% Adaptado do modelo de pôster de Nelson Lago (CCSL/IME-USP), licença MIT.
\documentclass[a1paper,brazilian]{article}
\usepackage{imegoodies}
\usepackage[poster,hidelinks,nobiblatex]{imelooks}
\pagestyle{empty}
\hypersetup{
  pdftitle={Testes ajudam a IA a gerar código mais correto?},
  pdfauthor={Eduardo Cruz Guedes},
}
\tcbposterset{fontsize=32pt}
\usetikzlibrary{arrows.meta,shapes.geometric}

\begin{document}
\begin{tcbposter}[poster={rows=12,columns=2,colspacing=1.2cm,rowspacing=.8cm}]

\posterbox[titlebox]{row=1,column=1,span=2,rowspan=2}{
  Testes ajudam a IA a gerar código mais correto?\par
  \medskip
  {\normalsize\sffamily\mdseries Eduardo Cruz Guedes \enspace\textbullet\enspace
  Bacharelado em Ciência da Computação, IME-USP\par
  Orientadores: Paulo Meirelles e Jorge Melegati}
}

\posterbox[adjusted title={Os dois pipelines}]{name=pipelines,row=3,column=1,span=2,rowspan=3.6}{
  \centering
  \begin{tikzpicture}[
      etapa/.style={draw=IMEframe,line width=2pt,fill=white,rounded corners=6pt,
        align=center,minimum height=3cm,text width=8.4cm,inner sep=8pt},
      gate/.style={draw=IMEframe,line width=2pt,fill=IMEtitlebg!15,diamond,
        aspect=1.6,inner sep=2pt,align=center,font=\bfseries},
      seta/.style={-{Stealth[length=18pt,width=16pt]},line width=3pt,IMEframe},
      grupo/.style={font=\bfseries,anchor=west},
    ]
    % Geração direta
    \node[grupo] at (-5,4.6) {Geração direta};
    \node[etapa] (b1) at (0,2.4) {Requisito};
    \node[etapa,text width=28cm] (b2) at (21.25,2.4) {Agente gera o código};
    \node[etapa] (b3) at (43.5,2.4) {Avaliação comum};
    \draw[seta] (b1) -- (b2);
    \draw[seta] (b2) -- (b3);
    % TDD
    \node[grupo] at (-5,-1.4) {TDD};
    \node[etapa] (t1) at (0,-3.6) {Requisito};
    \node[etapa] (t2) at (10.5,-3.6) {Agente escreve\\os testes};
    \node[etapa] (t3) at (21,-3.6) {Agente\\implementa};
    \node[gate] (t4) at (31,-3.6) {Testes\\passam?};
    \node[etapa] (t5) at (43.5,-3.6) {Avaliação comum};
    \draw[seta] (t1) -- (t2);
    \draw[seta] (t2) -- (t3);
    \draw[seta] (t3) -- (t4);
    \draw[seta] (t4) -- node[above] {sim} (t5);
    \draw[seta] (t4.south) -- ++(0,-1.4) -| node[pos=.25,below] {não}
      (t3.south);
  \end{tikzpicture}\par\bigskip
  \raggedright
  No TDD, um agente escreve os testes antes do código
  e o código volta para correção até passar neles.
}

\posterbox[adjusted title={Como comparamos}]{below=pipelines,column=1,span=2,rowspan=1.4}{
  TDD é desenvolvimento guiado por testes. Para medir a correção, usamos
  testes de avaliação comuns aos dois pipelines. Na avaliação de usabilidade,
  as notas da pessoa e do agente no navegador coincidiram em 5 de 30 comparações.
}

\posterbox[blankest,valign=top]{row=8,column=1,rowspan=3}{
  {\large\bfseries Formulários aprovados}\par
  Nos 22 critérios funcionais\par\bigskip
  \begin{tikzpicture}[x=8.5cm,y=2.9cm]
    \node[above,font=\bfseries] at (.5,2) {Geração direta};
    \node[above,font=\bfseries] at (1.5,2) {TDD};
    \node[left,font=\bfseries] at (-.03,1.5) {DeepSeek};
    \node[left,font=\bfseries] at (-.03,.5) {GLM-5.3};
    \foreach \coluna/\linha/\valor/\letra in {
      0/1/39/white,1/1/30/black,
      0/0/3/black,1/0/9/black%
    } {
      \pgfmathtruncatemacro{\intensidade}{2*\valor}
      \filldraw[fill=IMEtitlebg!\intensidade,draw=white,line width=1pt]
        (\coluna,\linha) rectangle ++(1,1);
      \node[font=\bfseries,text=\letra] at (\coluna+.5,\linha+.5) {\valor/50};
    }
  \end{tikzpicture}\par\bigskip
  \textbf{A abordagem com mais aprovações mudou conforme o modelo.}\par\medskip
  No DeepSeek, passaram 9 aplicações a mais com geração direta.
  No GLM-5.3, passaram 6 a mais com TDD.
}

\posterbox[blankest,valign=top]{row=8,column=2,rowspan=3}{
  {\large\bfseries Reconstrução do PreJanus}\par
  Aprovações em 85 testes reservados\par\bigskip
  \begin{tikzpicture}[x=.185cm,y=1.6cm]
    \foreach \total in {0,20,40,60,80} {
      \draw[gray!35] (\total,0) -- (\total,4.5);
      \node[below] at (\total,0) {\total};
    }
    \foreach \height/\valor/\shade/\variante in {
      4/62/IMEtitlebg/Geração direta,
      3/49/orange!80!brown/TDD com testes fixos,
      2/63/orange!80!brown/TDD com revisão,
      1/59/teal!70!black/Testes visíveis (rework)%
    } {
      \node[left] at (-1,\height) {\variante};
      \fill[\shade] (0,\height-.30) rectangle (\valor,\height+.30);
      \node[right,font=\bfseries] at (\valor,\height) {\valor};
    }
  \end{tikzpicture}\par\medskip
  Uma execução por variante na rodada inicial\par\medskip
  TDD com revisão custou cerca de 4 vezes mais que a geração direta.
}

\posterbox[adjusted title={Conclusão}]{row=11,column=1,span=2,rowspan=1.4}{
  A geração direta custou menos nos dois estudos. As abordagens com TDD
  consumiram mais tokens e tiveram ganhos de correção em algumas comparações.
}

\posterbox[footerbox]{column=1,span=2,above=bottom}{
  \href{https://educg550.github.io/tcc/}{Monografia, código e dados:
  educg550.github.io/tcc}
}

\end{tcbposter}
\end{document}
